% from ega4-2-fr.tex
\oldpage[IV-2]{221}

Nous supposerons encore que $\boldsymbol R$ vérifie la condition suivante :

$4^\circ$ Si $\boldsymbol P(Z,k)$ est vraie, alors, pour toute extension
\emph{de type fini} $k''$ de $k$, tous les anneaux locaux de
$Z\otimes_k k''$ vérifient la propriété $\boldsymbol R$.

On notera que ces conditions sont vérifiées lorsque l'on prend pour
$\mathcal C$ la catégorie des préschémas \emph{excellents} et pour
$\boldsymbol R$ une des propriétés (i) à (viii) de (7.5.3); pour la condition
$1^\circ$, cela résulte de (7.8.3, (iii)), et pour les conditions $2^\circ$ et
$3^\circ$, des raisonnements de (7.5.3), compte tenu de ce que tout anneau
excellent est caténaire. Enfin, pour la condition $4^\circ$, elle résulte,
pour (i), (ii) et (iii) de (6.7.1); pour (iv), (v), (vi) de (6.7.7), et pour
(viii) de (4.6.1); en ce qui concerne la propriété (vii) de (7.5.3), la
propriété correspondante $\boldsymbol P(Z,k)$ entraîne que $Z$ est localement
intègre (étant localement noethérien) et que chacun des sous-préschémas
intègres dont $Z$ est somme est géométriquement intègre, en vertu de (4.5.9)
et (4.6.1); on en conclut encore la propriété $4^\circ$ ci-dessus dans ce cas.

Avec ces notations et hypothèses :
\end{env}

\begin{proposition}[7.9.8]
\label{IV.2.s7.9.8-fr}
Soient $A$ un anneau local noethérien, $k$ son corps résiduel, $B$ un anneau
local noethérien tel que $\operatorname{Spec}(B)\in\mathcal C$,
$\varphi:A\to B$ un homomorphisme local faisant de $B$ un $A$-module plat;
posons $Y=\operatorname{Spec}(A)$, $X=\operatorname{Spec}(B)$ et soit
$f:X\to Y$ le morphisme correspondant à $\varphi$. On suppose que :

$1^\circ$ la propriété $\boldsymbol P(B\otimes_A k,k)$ soit vraie;

$2^\circ$ pour tout morphisme fini $Y_1\to Y$, on peut résoudre
$(Y_1)_{\mathrm{red}}$.

Alors, pour tout $y\in Y$, la propriété
$\boldsymbol P(f^{-1}(y),k(y))$ est vraie.
\end{proposition}

\begin{proof}
Notons que la condition $2^\circ$ de (7.9.8) est encore vérifiée quand on
remplace $Y$ par le spectre d'un anneau local en un idéal maximal d'une
$A$-algèbre finie $A'$ (7.9.2); d'autre part, en vertu de la condition
$1^\circ$ de (7.9.7), le spectre d'un anneau de fractions de
$B\otimes_AA'$ appartient aussi à $\mathcal C$. Le lemme (7.3.16.2) montre
alors (comme dans la partie I) du raisonnement de (7.5.2)) qu'il suffit de
prouver que, si $Y$ est intègre et si $y$ est le point générique de $Y$, les
anneaux locaux de la fibre $f^{-1}(y)$ vérifient $\boldsymbol R$.

Cela étant, il existe par hypothèse un préschéma régulier $Y'$ et un morphisme
propre et birationnel $g:Y'\to Y$; comme $Y'$ est noethérien et localement
intègre, l'hypothèse que $g$ est birationnel entraîne que $Y'$ est intègre;
soit $X'=X\times_YY'$, de sorte qu'on a un diagramme commutatif

\begin{env}[7.9.8.1]
\label{IV.2.s7.9.8.1-fr}
\[
\xymatrix{
X\ar[d]_f & X'\ar[l]_{g'}\ar[d]^{f'}\\
Y & Y'\ar[l]^g
}
\]
\end{env}

où $f'$ et $g'$ sont les projections canoniques. Compte tenu de la condition
$2^\circ$ de (7.9.7), le même raisonnement que dans le début de (7.9.3) montre
qu'il suffit de prouver

